\uuid{2lUL}
\titre{Matrice d'une application linéaire/Changement de bases/Puissance de matrices/Suites}
\auteur{hanani ; mzari}
\organisation{exo7}
\licence{CC BY-NC-SA 4.0}
\datecreate{2020-01-20}
\niveau{L1}
\module{Algèbre}
\chapitre{Matrice}
\sousChapitre{Matrice et application linéaire}
\difficulte{3}
\theme{Matrice d'une application linéaire, Changement de bases, Puissance de matrices, Suites}
\exo7qcmid{lille-2.matrice-application.029}

\contenu{
\begin{question}
\qclassification{108.03}
\qtags{motcle=Matrice d'une application linéaire/Changement de bases/Puissance de matrices/Suites}
Soit $f$ l'endomorphisme de $\Rr^4$ dont la matrice dans la base canonique  ${\cal B}$ est  :
$$A=\left(\begin{array}{rccc}1&0&0&0\\
0&1&-1&1\\ 0&-1&1&1\\0&0&0&-1\end{array}\right).$$
Soit ${\cal B}' = \{ a_1, a_2,  a_3, a_4\}$, où
$$a_1=(1,0,0,0),\; a_2=(0,1,1,0),\; a_3=(0,1,-1,0),\; a_4=(0,1,1,-1).$$
Soit $(u_n)_{n\ge 0}$, $(v_n)_{n\ge 0}$, $(w_n)_{n\ge 0}$ et $(k_n)_{n\ge 0}$ des suites récurrentes définies par la donnée
des réels $u_0, v_0,w_0,k_0$ et pour $n\ge 1$ :
$$(\mathtt{S})
\left\{\begin{array}{rcc}
u_n&=&u_{n-1}\\
v_n&=&v_{n-1}-w_{n-1}+k_{n-1}\\
w_n&=& -v_{n-1}+w_{n-1}+k_{n-1}\\
k_n&=&-k_{n-1}.\\
\end{array}\right.$$
Quelles sont les assertions vraies ?
\begin{answers}
\bad{$\{a_2\}$ est une base de $\ker (f-Id)$ et $\{a_1\}$ est une base de $\ker f$}.
\bad{$\{a_4 \}$ est une base de $\ker (f-2Id)$ et  $\{a_3\}$ est une base de $\ker (f+Id)$}.
\good{${\cal B}'$ est une base de $\Rr^4$ et la matrice de $f$ dans cette base est :
$$B=\left(\begin{array}{rccc}1&0&0&0\\0&0&0&0\\
0&0&2&0\\0&0&0&-1\end{array}\right).$$}
\good{Pour tout entier $n\ge 1$, on a : $$(\mathtt{S})
\left\{\begin{array}{rcc}
u_n&=&u_0\\
v_n&=&2^{n-1}v_0-2^{n-1}w_0+(-1)^{n-1}k_0\\
w_n&=&-2^{n-1}v_0+2^{n-1}w_0+(-1)^{n-1}k_0 \\
k_n&=&(-1)^nk_0.\\
\end{array}\right.$$ }
\end{answers}
\begin{explanations} On vérifie que
$\{a_1\}$ est une base de $\ker (f-Id)$,  $\{a_2\}$ est une base de $\ker f$, $\{a_3 \}$ est une base de $\ker (f-2Id)$,   $\{a_4\}$ est une base de $\ker (f+Id)$ et que ${\cal B}'$ est une base de $\Rr^4$.
\vskip0mm
La matrice de $f$ dans la base ${\cal B}'$ est :
$B= \left(\begin{array}{rccc}
1&0&0&0\\
0&0&0&0\\
0&0&2&0\\
0&0&0&-1\\
\end{array}\right)$.
\vskip0mm
La matrice de passage de  ${\cal B}$ à   ${\cal B}'$ (c.à.d la matrice de l'identité de $\Rr_3[X]$ de
la base ${\cal B}'$ à la base ${\cal B}$)
est : $P= \left(\begin{array}{rccc} 1&0&0&0\\
0&1&1&1\\
0&1&-1&1\\ 0&0&0&-1\\
\end{array}\right)\Rightarrow \displaystyle P^{-1}= \frac{1}{2}\left(\begin{array}{rccc}
2&0&0&0\\
0&1&1&2\\
0&1&-1&0\\
0&0&0&-2\\
\end{array}\right).$
\vskip0mm
De la relation : $A=PBP^{-1}$, on déduit que
$A^n= PB^nP^{-1}=\left(\begin{array}{rccc}
1&0&0&0\\
0 &2^{n-1}&-2^{n-1}&(-1)^{n-1}\\
0 &-2^{n-1}&2^{n-1}&(-1)^{n-1}\\
0 &0&0&(-1)^n\\
\end{array}\right)$, pour tout entier $n\ge 1$.\\
De la relation : $\left(\begin{array}{r}u_n\\v_n\\ w_n\\k_n\end{array}\right) = A  \left(\begin{array}{r}
u_{n-1}\\v_{n-1}\\ w_{n-1}\\k_{n-1}\end{array}\right)$, on déduit que :
$\left(\begin{array}{r}u_n\\v_n\\ w_n\\k_n\end{array}\right) = A^n \left(\begin{array}{r}u_0\\v_0\\ w_0\\k_0\end{array}\right)$.
\end{explanations}
\end{question}
}
